\section{Formalizaci\'on del requisito}
El handbook presenta la formalizaci\'on como un proceso iterativo que va desde una necesidad impl\'icita hasta un requisito expl\'icito, entendible, revisable y gestionable. Una secuencia pr\'actica de pasos es la siguiente:

\begin{enumerate}[leftmargin=1.2em]
  \item \textbf{Identificar la necesidad.} El punto de partida es una necesidad, deseo o problema percibido por un stakeholder \cite[p.~10]{cpre}.
  \item \textbf{Elicitar la informaci\'on.} Hay que extraer las necesidades ocultas en sus fuentes y convertirlas en requisitos expl\'icitos usando t\'ecnicas de recopilaci\'on, colaboraci\'on, observaci\'on o an\'alisis de artefactos \cite[pp.~88--95]{cpre}.
  \item \textbf{Explorar ideas y alternativas.} Cuando el problema admite soluciones innovadoras, se usan t\'ecnicas de dise\~no e ideaci\'on como brainstorming, prototipos o escenarios \cite[pp.~96--99]{cpre}.
  \item \textbf{Resolver conflictos.} Las necesidades de diferentes stakeholders pueden chocar; antes de avanzar, deben identificarse, analizarse y resolverse mediante consenso, compromiso, votaci\'on, overruling o definici\'on de variantes \cite[pp.~100--108]{cpre}.
  \item \textbf{Documentar con la representaci\'on adecuada.} El requisito puede expresarse como texto natural, plantilla o modelo. El handbook explica que un work product puede ir desde una simple frase hasta una especificaci\'on completa de requisitos, y que la forma elegida depende del nivel de abstracci\'on y del prop\'osito \cite[pp.~30--32]{cpre}.
  \item \textbf{Validar.} Antes de entregar el requisito, se debe comprobar que refleja las necesidades reales y que los stakeholders est\'an de acuerdo; la validaci\'on debe empezar dentro de RE, no al final del proyecto \cite[pp.~109--111]{cpre}.
  \item \textbf{Gestionar ciclo de vida, versiones y cambios.} Una vez formalizado, el requisito debe quedar bajo control de versiones, trazabilidad y gesti\'on del cambio para mantener coherencia entre estados, configuraciones y baselines \cite[pp.~131--140]{cpre}.
\end{enumerate}
